\section{Discussion}\label{sec:discussion}

We bring together the functional and Hamiltonian correspondences to
explain the common BA setting and identify directions for further work.

The canonical multiwave system provides a canonical
model for scalar BA dynamics whose interaction structure is independent
of the curve type. The two-way correspondence for arbitrary positive $L,M,N$ and the
higher-flow result give this interpretation its substance: prescribed
BA evaluations find a home in the cubic equations, generic canonical
initial modes recover their BA data, and the spectral Hamiltonians
supply every Abel direction on birational spectral levels together
with the mode phases.
This correspondence motivates further study of the system's
Hamiltonian hierarchy and reductions across different curve families.
The invariant real forms in \cref{sec:canonical-modes} apply to real
initial data; explicit curve-and-divisor reality criteria, particularly
for decay-form BA realizations, remain to be developed.

The realization criterion turns prescribed curve data into canonical
multiwave dynamics: residue equations determine compatible weights,
and opposite-group interpolation values determine the linear phase
rates. The kernel lemma gives the differential identities behind this
construction. The matrix constrained-KP correspondence in
\cref{sec:BA-constrained-KP} packages these identities into two
commuting finite-rank operators for arbitrary $N$. Choosing the local
spectral coordinate fixes their KP times and source evolution;
spectral Hamiltonians and power charges give their canonical
generators on each birational BA level. The one-point member recovers
scalar operators and their conserved coefficients.
\Cref{prop:BA-prescribed-spectral-curve,thm:BA-spectral-converse}
identify the realized curve and reconstruct BA data from generic
canonical initial modes of the full three-array system.

The comparison with Adams, Harnad and Hurtubise gives two compatible
views of the same matrix evolution. Their canonical moment-map lift
has quartic interactions, while the cubic modes use the constant
member of the matrix Poisson pencil. Complete-fiber BA convolution
turns the bilinear lift into endpoint BA functions and expresses this
change of interaction order as a Sylvester displacement relation.
This connection complements the prescribed-point realization: it
explains how the direct cubic modes sit within the established
spectral construction while retaining a simple canonical bracket.

For the full three-matrix family, the spectral and principal-moment
integrals give the complete Liouville count for every fixed phase
profile, with the detailed proof in Appendix~\ref{sec:integrability-proof}. The symmetry interpretation in
\cref{sec:gauge-Abel-symmetries} distinguishes mode gauges, equal-phase
changes of basis and flows that become translations in spectral
divisor coordinates. These different symmetries explain why power
charges alone need not complete the integral count.

In the one-point hyperelliptic specialization, the Mumford projection
has a useful optical analogy in the passage from
complex field modes to Stokes variables. In elliptically solvable optical
systems, Stokes variables describe closed projected dynamics and determine
the logarithmic rates needed to reconstruct the individual modes.
The Mumford variables play a similar role here: the powers and their
derivatives determine the moving divisor polynomials, whose evolution
closes and whose values determine the BA logarithmic rates in
\cref{eq:BA-powers-rates}. The BA formulas give the mode reconstruction
explicitly, with the remaining gauge choices carried by the constant
normalizations $c_k$. Thus both descriptions connect projected dynamics
to the evolution of the underlying complex modes. The Stokes projection
omits one common optical phase, while the Mumford projection here omits
the $K$ independent reciprocal mode gauges.

The broader spectral family offers further function-theoretic problems.
Its marked plane models describe different covering degrees and branches
at infinity; special levels and involutions produce lower-genus
normalizations and Prym descriptions. The BA construction applies across these curve types: one puncture
gives the vector--matrix--vector member, while several punctures
produce the full arrays through a common scalar BA basis and a
rank-$N$ kernel derivative. The puncture count is independent of
the genus or the number of infinities in another plane presentation.
Singular spectral levels and nonsemisimple moment
blocks also require limiting constructions beyond the generic converse.

The commuting spectral Hamiltonians give a multi-time view of the
interactions. \Cref{thm:BA-higher-flows} computes their puncture times and constant
mode phases by finite meromorphic interpolation. On a birational
level, finitely many spectral coefficients generate every Abel direction.
A scalar evolution selects one such direction, and higher Hamiltonians
introduce higher-order mixing terms on the same mode space. This makes
the translation and gauge symmetries explicit in the BA representation.
The interaction orders and combinations associated with particular
directions provide a further problem for constructing polynomial
Hamiltonians from BA identities.
